\section{The algebra of time-invariant observables}
\label{ex-phases}

The Hilbert space constructed from an unselected correlation table need not
have a one-dimensional vacuum space. We first identify the whole subspace
left fixed by time evolution. Compressed slab observables on this subspace
will then distinguish its phases. Throughout this section the hypotheses
of Theorem~\ref{ex-euclidean-semigroup} remain in force. We write
\(\mathcal H=\mathcal H_E\) and \(\Omega=\Omega_E\). In this section \(e_0\) denotes the chosen time axis and \(x_0\) its coordinate; this local notation does not change the earlier coordinate ordering.

\begin{proposition}[The large-time projection]\label{ex-ground-projection}
The subspace
\[
 \mathcal H_0=\{v\in\mathcal H:E_tv=v\text{ for every }t\ge0\}
\]
is closed. If \(P_0\) denotes its orthogonal projection, then
\(E_tv\to P_0v\) in norm as \(t\to\infty\), for every \(v\in\mathcal H\).
\end{proposition}
\begin{proof}
Each \(E_t-I\) is continuous, so its kernel is closed; their intersection
is \(\mathcal H_0\). Fix \(v\). Contractivity and the semigroup law show
that \(a(t)=\|E_tv\|^2\) decreases to its nonnegative infimum \(\ell\).
Self-adjointness gives, for \(s\le t\),
\[
 \langle E_sv,E_tv\rangle
 =\langle v,E_{s+t}v\rangle
 =\|E_{(s+t)/2}v\|^2.
\]
The scalar on the right is real. Consequently
\[
 \|E_sv-E_tv\|^2
 =a(s)+a(t)-2a((s+t)/2)\le a(s)-a(t)\le a(s)-\ell.
\]
For large \(s,t\) this tends to zero, after exchanging \(s,t\) if
necessary. Completeness gives a norm limit \(v_0\). For fixed \(r\ge0\),
continuity of \(E_r\) and \(E_rE_t=E_{r+t}\) give \(E_rv_0=v_0\).
For \(z\in\mathcal H_0\), self-adjointness gives
\(\langle E_tv,z\rangle=\langle v,E_tz\rangle=\langle v,z\rangle\).
Passing to the limit shows \(v-v_0\perp\mathcal H_0\). These two
properties characterize the orthogonal projection, so \(v_0=P_0v\).
\end{proof}

We next explain the analytic estimate needed to insert an entire word as a
bounded operator. A word \(w\) lies in a slab of width \(d>0\) if all its
points are distinct and have time coordinates in \((0,d)\). For a positive
word \(v\), put
\[
 I_wv=w\,\tau_dv,\qquad w^\sharp=\tau_d\theta w.
\]
Both are positive words, and \(w^\sharp\) lies in the same slab. Let
\(I_w\) also denote the linear map on finite formal word sums. Translation
invariance and reflection give
\begin{equation}\label{ex-slab-formal-adjoint}
 \langle[I_wc],[b]\rangle=\langle[c],[I_{w^\sharp}b]\rangle.
\end{equation}
Indeed, translating the configuration
\(\theta(w\tau_dv)\,u=(\tau_{-d}\theta v)(\theta w)u\)
by \(d\) gives \((\theta v)(\tau_d\theta w)(\tau_du)\), the
configuration in the right-hand pairing. This proves the identity on
words and then on finite sums.

\begin{lemma}[From repeated insertion to a bounded operator]
\label{ex-train-bound}
Fix a slab word \(w\) and \(D>0\). Suppose that for every positive word
\(v\) there is \(B_v\ge0\) such that
\begin{equation}\label{ex-train-input}
 \bigl\|[(I_{w^\sharp}I_w)^nv]\bigr\|^2\le B_vD^{4n}
 \quad(n=0,1,2,\ldots).
\end{equation}
There is a unique bounded operator \(A_d(w)\) on \(\mathcal H\) with
\(A_d(w)[v]=[w\tau_dv]\); it satisfies \(\|A_d(w)\|\le D\).
\end{lemma}
\begin{proof}
Write \(T=I_{w^\sharp}I_w\) on formal sums and \(Xc=[c]\).
The triangle inequality and \eqref{ex-train-input} imply, for
\(c=\sum_vc_vv\),
\[
 \|XT^nc\|\le C_cD^{2n},\qquad C_c=\sum_v|c_v|\sqrt{B_v}<\infty.
\]
Equation~\eqref{ex-slab-formal-adjoint} shows that \(T\) is symmetric
for the possibly degenerate form induced by \(X\). Repeatedly moving
\(T\) across the inner product therefore gives
\(\|XT^nc\|^2=\langle Xc,XT^{2n}c\rangle\).
Put \(b_n=D^{-2n}\|XT^nc\|\) and \(M=\sup_{n\ge0}b_n<\infty\).
Cauchy--Schwarz gives \(b_n^2\le\|Xc\|b_{2n}\le\|Xc\|M\).
Taking the supremum yields \(M^2\le\|Xc\|M\), hence \(M\le\|Xc\|\).
Finally,
\[
 \|XI_wc\|^2=|\langle Xc,XTc\rangle|
 \le\|Xc\|\|XTc\|\le D^2\|Xc\|^2.
\]
Thus insertion respects the null space and induces a bounded operator on
the quotient. Its unique extension to the completion has the asserted
norm and action. Uniqueness follows from density of the word span.
\end{proof}

The estimate in \eqref{ex-train-input} concerns arbitrarily long repeated
words; the constants \(B_v\) may depend on the exterior word. Bounds for
each fixed correlation order alone do not supply this estimate. Its
production from the lattice construction remains an input at this stage.
Assume henceforth that it holds, with some \(D\), for every slab word.
The resulting operators obey
\begin{align}
 A_d(\varnothing)&=E_d,\label{ex-slab-laws}\\
 A_{d+e}(w\tau_dv)&=A_d(w)A_e(v),\notag\\
 A_d(w)^*&=A_d(w^\sharp),\notag\\
 A_{d+r}(w)&=A_d(w)E_r,\qquad
 A_{r+d}(\tau_rw)=E_rA_d(w).\notag
\end{align}
Here each word lies in its indicated slab and \(r\ge0\).
The first, second and last two identities follow by applying both sides
to a positive word and concatenating and translating its entries.
The adjoint identity follows from \eqref{ex-slab-formal-adjoint}.
All extend from the dense word span by boundedness. Spatial translation
also gives
\(U_yA_d(w)=A_d(w+y)U_y\), by the same calculation on words.

\subsection{Spatial invariance of the ground space}
Assume now full orthogonal covariance of the correlation table: there is
an action on quotient field labels such that, for every \(K\in O(4)\) and every
distinct-point word \(w\), the jointly rotated word \(Kw\) satisfies
\(S(Kw)=S(w)\). Rotation includes this label action in what follows. To evaluate a rotated word, project each raw label to its field class, rotate that class, and choose a raw representative. Correlations are independent of that choice. No group action on the chosen raw representatives is assumed.
This is stronger than covariance under signed coordinate permutations.

\begin{lemma}\label{ex-ground-spatial}
Every \(v\in\mathcal H_0\) is fixed by all spatial translations.
\end{lemma}
\begin{proof}
For two distinct-point words \(p,q\), set
\(F_{p,q}(x)=S(p(q+x))\) whenever the two clusters are disjoint.
Translate \(\theta p\) and \(q\) separately far enough in positive
time to obtain positive words \(v,w\). If the translation amounts are
\(r,s\), then, whenever \(x_0\ge r+s\), translation invariance gives
\begin{equation}\label{ex-separated-pair}
 F_{p,q}(x)=\langle[v],U_{x_\perp}E_{x_0-r-s}[w]\rangle,
 \qquad x_\perp=x-x_0e_0.
\end{equation}
The time separation also ensures that the two clusters are disjoint.
For any fixed \(h\in\mathbb R\) and any path \(x(t)\) with
\(x_0(t)\to\infty\), it follows that
\begin{equation}\label{ex-normal-shift}
 F_{p,q}(x(t)+he_0)-F_{p,q}(x(t))\longrightarrow0.
\end{equation}
Indeed, the difference of the two heat-evolved vectors in
\eqref{ex-separated-pair} tends to zero by
Proposition~\ref{ex-ground-projection}; the spatial unitary has norm one,
even when \(x_\perp(t)\) varies without bound.

If \(b_0>0\), choose an orthogonal map \(K\) taking
\(b\) to \(|b|e_0\). Orthogonality implies
\((Ke_0)_0=b_0/|b|>0\). Rotate both clusters and apply
\eqref{ex-normal-shift} along the path \(tKe_0\), whose time
coordinate tends to infinity. Covariance gives
\[
 F_{p,q}(te_0+b)-F_{p,q}(te_0)\longrightarrow0.
\]
For a spatial \(y\), take \(b=y+e_0\) and subtract
\eqref{ex-normal-shift} along \(x(t)=te_0+y\), with \(h=1\).
We obtain \(F_{p,q}(te_0+y)-F_{p,q}(te_0)\to0\).
Apply this with \(p=\theta v,q=w\) for positive words. The two
correlations are \(\langle[v],U_yE_t[w]\rangle\) and
\(\langle[v],E_t[w]\rangle\). Their limits show
\(\langle[v],(U_y-I)P_0[w]\rangle=0\).
Density of word vectors first gives \((U_y-I)P_0[w]=0\) and then
\((U_y-I)P_0=0\). Since \(P_0\) is the identity on \(\mathcal H_0\),
the conclusion follows.
\end{proof}

\subsection{Compression and separation of clusters}
For a slab word define
\(M(w)=P_0A_d(w)|_{\mathcal H_0}\). The width is immaterial:
increasing it appends an \(E_r\), which is the identity on
\(\mathcal H_0\); two widths can be compared with their sum.
Likewise a positive time shift inserts \(E_r\) on the left, and
\(P_0E_r=P_0\). Spatial shifts give the same operator by
Lemma~\ref{ex-ground-spatial} and spatial covariance. For a negative
time shift, reverse the comparison. Every finite distinct-point word can
be translated into a slab. Thus \(M(w)\) is well-defined for every such
word, independent of both the placement and width, and
\begin{equation}\label{ex-compression-values}
 M(w+b)=M(w),\qquad
 \langle\Omega,M(w)\Omega\rangle=S(w).
\end{equation}
The latter follows from \(A_d(w)\Omega=[w]\) after a slab placement,
\(P_0\Omega=\Omega\), and translation invariance.

For a positive word also put \(X(w)=A_d(w)|_{\mathcal H_0}\), as a map
into \(\mathcal H\). Width independence follows as above, and
\[
 X(w)\Omega=[w],\qquad M(w)\Omega=P_0[w],\qquad
 M((\theta v)w)=X(v)^*X(w).
\]
For the last identity, place the reflected word into the slab of width
\(d\), append the word \(w\) in a subsequent slab of width \(e\),
and use the concatenation and adjoint identities in
\eqref{ex-slab-laws}. Compression then gives exactly the displayed
product. In particular \(M(\theta v)=M(v)^*\). This also shows that
the generators are closed under adjoints, after translation into slabs.

\begin{lemma}[Separated-cluster limit]\label{ex-cluster-train}
Let \(w_1,\ldots,w_m\) be words with distinct points. Suppose
\(b_1(n),\ldots,b_m(n)\in\mathbb R^4\) satisfy
\((b_{j+1}(n)-b_j(n))_0\to\infty\) for \(j<m\). Then
\[
 S((w_1+b_1(n))\cdots(w_m+b_m(n)))
 \longrightarrow\langle\Omega,M(w_1)\cdots M(w_m)\Omega\rangle.
\]
The concatenated words have distinct points for all sufficiently large \(n\).
\end{lemma}
\begin{proof}
Fixed translations of individual words, compensated in their \(b_j\),
preserve the assertion and the gap condition. Make all words positive.
We prove a stronger vector limit by induction from the last cluster:
after translating the concatenation so that its first cluster has its
original position, its vector tends to
\(X(w_1)M(w_2)\cdots M(w_m)\Omega\).
For one cluster this is \(X(w_1)\Omega=[w_1]\).

Choose a slab width \(d\) for \(w_1\). Write the displacement of the
second cluster relative to the first as \(g_ne_0+y_n\), with \(y_n\)
spatial. Eventually \(g_n>d\). If \(z_n\) is the positive word vector
of the remaining clusters relative to the second, the whole vector is
\(A_d(w_1)U_{y_n}E_{g_n-d}z_n\).
By induction \(z_n\to z=X(w_2)M(w_3)\cdots M(w_m)\Omega\).
The estimates
\[
 \|U_{y_n}E_{g_n-d}z_n-P_0z\|
 \le\|z_n-z\|+\|E_{g_n-d}z-P_0z\|\longrightarrow0
\]
use contraction, unitarity and \(U_{y_n}P_0z=P_0z\). Boundedness of
\(A_d(w_1)\) proves the vector limit. The same induction shows positivity
and distinctness eventually, since the successive time gaps tend to
infinity. Pairing the vector limit with \(\Omega\), and using translation
invariance and \(P_0\Omega=\Omega\), proves the scalar assertion.
The empty list has correlation one and needs no induction.
\end{proof}

\begin{proposition}[Commutativity]\label{ex-ground-commutative}
The norm-closed unital star algebra \(\mathcal A\subset B(\mathcal H_0)\)
generated by the \(M(w)\) is commutative, and \(\Omega\) is cyclic and
separating for \(\mathcal A\).
\end{proposition}
\begin{proof}
Write \(F(w_1,\ldots,w_m)=\langle\Omega,\prod_jM(w_j)\Omega\rangle\),
in the indicated order. If \(b_0>0\) and \((Kb)_0>0\), place the
clusters at \(jnb\), rotate the configuration by \(K\), and use
Lemma~\ref{ex-cluster-train} before and after rotation. Covariance gives
\(F(w_1,\ldots,w_m)=F(Kw_1,\ldots,Kw_m)\).

To exchange the middle two of four words, use one spatial axis \(e_1\)
and displace the four clusters by
\[
 0,\quad n(e_0+2e_1),\quad n(2e_0+e_1),\quad3n(e_0+e_1).
\]
They are ordered by time as written. Interchange \(e_0,e_1\) by an
orthogonal map \(K\). The resulting time order is first, third, second,
fourth. Permutation symmetry of distinct-point correlations and the
cluster limit give
\(F(x,y,w,z)=F(Kx,Kw,Ky,Kz)\).
Use the preceding rotation identity with \(b=e_0+e_1\), fixed by \(K\),
to conclude \(F(x,y,w,z)=F(x,w,y,z)\).

The vectors \(M(v)\Omega=P_0[v]\), for positive words \(v\), have
dense linear span in \(\mathcal H_0\): approximate any vector in
\(\mathcal H_0\) by word sums in \(\mathcal H\) and apply \(P_0\).
For two generators, their commutator has zero matrix elements between
these dense vectors, since
\[
 \langle M(u)\Omega,[M(w),M(z)]M(v)\Omega\rangle
 =F(\theta u,w,z,v)-F(\theta u,z,w,v)=0.
\]
The commutator is bounded, hence is zero. The generators and their
adjoints therefore commute; polynomial approximation and continuity of
multiplication show that their norm closure is commutative. The same
density proves cyclicity. If \(A\Omega=0\), then
\(AB\Omega=BA\Omega=0\) for every \(B\in\mathcal A\).
Cyclicity and boundedness give \(A=0\), which is separation.
\end{proof}

This proves commutativity from the displayed slab and covariance inputs.
The next construction must establish norm separability of \(\mathcal A\),
represent the vacuum state on its compact character space, and show that
each resulting whole-word correlation table retains the required
regularity, null spaces and field identities. Those conclusions are not
part of Proposition~\ref{ex-ground-commutative}.
